% © 2026 Ing. David Jaroš — CC BY-NC-ND 4.0
%
% This work is licensed under a Creative Commons Attribution-NonCommercial-NoDerivatives
% 4.0 International License (CC BY-NC-ND 4.0).
%
% License History: Earlier drafts (up to v0.3) were released under CC BY 4.0.
% From v0.4 onward, all material is released under CC BY-NC-ND 4.0 to protect
% the integrity of the theoretical work during ongoing academic development.

\documentclass[12pt]{article}
\usepackage{amsmath,amssymb,amsthm}
\usepackage{geometry}
\usepackage{hyperref}
\usepackage{tcolorbox}
\usepackage{xcolor}
\geometry{margin=1in}

% Theorem-like environments
\newtheorem{definition}{Definition}
\newtheorem{proposition}{Proposition}
\newtheorem{theorem}{Theorem}
\newtheorem{remark}{Remark}
\newtheorem{conjecture}{Conjecture}

\title{ST-2: Interaction Between Real and Imaginary Sectors in UBT\\
Cross-Terms in the Action and Energy-Momentum\\[0.5em]
\large\textsc{Speculative Extension — Invisibility Track}}
\author{David Jaroš}
\date{2026-03-03}

\begin{document}
\maketitle

\begin{tcolorbox}[colback=yellow!10, title=Status]
\textbf{SPECULATIVE EXTENSION} — This content is not part of
the canonical UBT framework.  It is an exploratory investigation
and has not been validated against the canonical field equations.
\end{tcolorbox}

\vspace{1em}
\noindent\textit{Output of Sub-task ST-2 from the Invisibility / Imaginary Sector
investigation track.  This is the \emph{key} sub-task: if cross-terms in
$S[\Theta_\mathrm{real} + \Theta_\mathrm{imag}]$ vanish identically, the
sectors are fully decoupled; otherwise partial (gravitational) invisibility
is possible.  See also:
\texttt{st1\_imaginary\_object\_definition.tex},
\texttt{st3\_stability.tex},
\texttt{st4\_ctc\_invisibility\_connection.tex},
\texttt{st5\_falsifiability.md}.}

%----------------------------------------------------------------------
\section*{Abstract}
%----------------------------------------------------------------------
We compute the cross-terms in the UBT action
$S[\Theta_\mathrm{real} + \Theta_\mathrm{imag}]$ at leading order.
Two cases are distinguished: (A) a metric bilinear in $\Theta$,
for which cross-terms generically do \emph{not} vanish, leading to
\emph{partial invisibility} (gravitationally detectable but
electromagnetically invisible); and (B) special configurations
where cross-terms cancel, yielding exact sector decoupling.
The energy-momentum tensor $T_{\mu\nu}[\Theta_\mathrm{imag}]$ is
evaluated and shown to be real (under mild assumptions), implying
that imaginary-sector objects can bend light.  The phenomenology
is analogous to gravitationally interacting dark matter.
All conclusions are \textbf{speculative}.

%----------------------------------------------------------------------
\section{Setup: Total Field and Action Decomposition}
%----------------------------------------------------------------------

We split the total $\Theta$-field as
\begin{equation}
  \Theta \;=\; \Theta_R + \Theta_I,
  \label{eq:field_split}
\end{equation}
where $\Theta_R$ (real-sector) and $\Theta_I$ (imaginary-sector)
are defined via
\begin{equation}
  \Theta_R \;:=\; \operatorname{Re}(\Theta_0)\,\mathbf{1}
            + \sum_{k=1}^{3}\operatorname{Re}(\Theta_k)\,\mathbf{e}_k,
  \qquad
  \Theta_I \;:=\; i\operatorname{Im}(\Theta_0)\,\mathbf{1}
            + \sum_{k=1}^{3}i\operatorname{Im}(\Theta_k)\,\mathbf{e}_k,
\end{equation}
with $\{\mathbf{1}, \mathbf{e}_1, \mathbf{e}_2, \mathbf{e}_3\}$ the
standard quaternion basis.  By construction,
$\Theta_R^* = \Theta_R$ (no complex phase) and
$\Theta_I = i\,\widetilde\Theta_I$ for some purely real biquaternion
$\widetilde\Theta_I$.

The UBT action (schematically) reads
\begin{equation}
  S[\Theta, \mathcal{G}]
  \;=\; \int d^4x\sqrt{-g}\left[
    \tfrac{1}{2\kappa}R
    - \partial_\mu\Theta^\dagger \partial^\mu\Theta
    - V(\Theta^\dagger\Theta)
  \right],
  \label{eq:action_gen}
\end{equation}
where $g = \det(g_{\mu\nu})$, $R$ is the Ricci scalar of $g_{\mu\nu}$,
and indices are raised with $g^{\mu\nu}$.

%----------------------------------------------------------------------
\section{Cross-Terms in the Kinetic Sector}
%----------------------------------------------------------------------

Inserting $\Theta = \Theta_R + \Theta_I$ into the kinetic term:
\begin{align}
  \partial_\mu\Theta^\dagger \partial^\mu\Theta
  &\;=\; \partial_\mu(\Theta_R + \Theta_I)^\dagger
         \partial^\mu(\Theta_R + \Theta_I) \nonumber\\
  &\;=\; \partial_\mu\Theta_R^\dagger \partial^\mu\Theta_R
       + \partial_\mu\Theta_I^\dagger \partial^\mu\Theta_I
       + \underbrace{\partial_\mu\Theta_R^\dagger \partial^\mu\Theta_I
       + \partial_\mu\Theta_I^\dagger \partial^\mu\Theta_R}_{
         \text{cross-terms } \mathcal{X}_\mathrm{kin}}.
  \label{eq:kinetic_expand}
\end{align}

\paragraph{Evaluation of $\mathcal{X}_\mathrm{kin}$.}
Let $\Theta_I = i\,\widetilde\Theta_I$ with $\widetilde\Theta_I$
real (meaning all quaternion components are real-valued).  Then:
\begin{align}
  \partial_\mu\Theta_R^\dagger \partial^\mu\Theta_I
  &\;=\; \partial_\mu\Theta_R^\dagger \,(i\,\partial^\mu\widetilde\Theta_I)
  \;=\; i\,\partial_\mu\Theta_R^\dagger \,\partial^\mu\widetilde\Theta_I, \\
  \partial_\mu\Theta_I^\dagger \partial^\mu\Theta_R
  &\;=\; (-i)\,\partial_\mu\widetilde\Theta_I^\dagger\,\partial^\mu\Theta_R.
\end{align}
The sum is
\begin{equation}
  \mathcal{X}_\mathrm{kin}
  \;=\; i\,\partial_\mu\Theta_R^\dagger \partial^\mu\widetilde\Theta_I
       - i\,\partial_\mu\widetilde\Theta_I^\dagger \partial^\mu\Theta_R
  \;=\; 2i\,\operatorname{Im}\!\left(
        \partial_\mu\Theta_R^\dagger \partial^\mu\widetilde\Theta_I
        \right).
  \label{eq:cross_kin_result}
\end{equation}
This is \emph{purely imaginary} (as a complex number).  When the
action integral is taken over real spacetime, $\operatorname{Re}(S)$
receives no contribution from $\mathcal{X}_\mathrm{kin}$ at this order.

\begin{proposition}[Kinetic cross-term is purely imaginary]
\label{prop:kin_cross}
At the level of the kinetic action \eqref{eq:kinetic_expand}, the
cross-term $\mathcal{X}_\mathrm{kin}$ is purely imaginary and does
not contribute to $\operatorname{Re}(S)$.  Consequently, the real
part of the action decouples at kinetic order:
\begin{equation}
  \operatorname{Re}(S_\mathrm{kin})
  \;=\; \operatorname{Re}(S_\mathrm{kin}[\Theta_R])
       + \operatorname{Re}(S_\mathrm{kin}[\Theta_I]).
\end{equation}
\end{proposition}

%----------------------------------------------------------------------
\section{Cross-Terms in the Potential Sector}
%----------------------------------------------------------------------

The potential $V(\Theta^\dagger\Theta)$ expanded around
$\Theta = \Theta_R + \Theta_I$:
\begin{equation}
  V((\Theta_R+\Theta_I)^\dagger(\Theta_R+\Theta_I))
  \;=\; V(\Theta_R^\dagger\Theta_R
         + \Theta_R^\dagger\Theta_I
         + \Theta_I^\dagger\Theta_R
         + \Theta_I^\dagger\Theta_I).
\end{equation}
The cross terms in the argument are:
\begin{align}
  \Theta_R^\dagger\Theta_I + \Theta_I^\dagger\Theta_R
  &\;=\; \Theta_R^\dagger(i\widetilde\Theta_I)
        + (-i\widetilde\Theta_I^\dagger)\Theta_R \nonumber\\
  &\;=\; i(\Theta_R^\dagger\widetilde\Theta_I
          - \widetilde\Theta_I^\dagger\Theta_R)
  \;=\; 2i\,\operatorname{Im}(\Theta_R^\dagger\widetilde\Theta_I).
  \label{eq:pot_cross_arg}
\end{align}
The argument of $V$ acquires a purely imaginary additive term.
Taylor-expanding $V$ to first order in this imaginary shift:
\begin{equation}
  V(\Theta_R^\dagger\Theta_R + 2i\,\delta)
  \;\approx\; V(\Theta_R^\dagger\Theta_R)
             + 2i\,\delta\,V'(\Theta_R^\dagger\Theta_R)
             + O(\delta^2),
\end{equation}
where $\delta = \operatorname{Im}(\Theta_R^\dagger\widetilde\Theta_I)$
is a real quantity.  This first-order correction is purely imaginary
and therefore does not contribute to $\operatorname{Re}(S)$.
The second-order term $V'' \delta^2$ is \emph{real} and couples
the sectors.

\begin{proposition}[Potential cross-terms at second order]
\label{prop:pot_cross}
Cross-terms between $\Theta_R$ and $\Theta_I$ in the potential
$V(\Theta^\dagger\Theta)$ first appear at second order in $\Theta_I$:
\begin{equation}
  \delta_\mathrm{cross}V
  \;\sim\; -2\,\left[
    \operatorname{Im}(\Theta_R^\dagger\widetilde\Theta_I)
  \right]^2 V''(\Theta_R^\dagger\Theta_R)
  + O(\Theta_I^3).
\end{equation}
This term is real and contributes to the equations of motion at
order $|\Theta_I|^2$.
\end{proposition}

%----------------------------------------------------------------------
\section{Effect on the Effective Einstein Equations}
%----------------------------------------------------------------------

The effective stress-energy sourcing the real Einstein equations is
\begin{equation}
  T_{\mu\nu}^{\mathrm{eff}}
  \;=\; \operatorname{Re}\!\left(\mathcal{T}_{\mu\nu}[\Theta,\mathcal{G}]\right)
  \;=\; T_{\mu\nu}[\Theta_R]
       + T_{\mu\nu}[\Theta_I]
       + T_{\mu\nu}^{\mathrm{cross}}[\Theta_R,\Theta_I].
\end{equation}

\paragraph{Individual contributions.}
From the kinetic piece:
\begin{align}
  T_{\mu\nu}[\Theta_R]
  &\;=\; \partial_\mu\Theta_R^\dagger\partial_\nu\Theta_R
        - \tfrac{1}{2}g_{\mu\nu}(\partial\Theta_R)^2
        - \tfrac{1}{2}g_{\mu\nu}V(\Theta_R^\dagger\Theta_R),\\
  T_{\mu\nu}[\Theta_I]
  &\;=\; \partial_\mu\Theta_I^\dagger\partial_\nu\Theta_I
        - \tfrac{1}{2}g_{\mu\nu}(\partial\Theta_I)^2
        - \tfrac{1}{2}g_{\mu\nu}V(\Theta_I^\dagger\Theta_I).
\end{align}

\paragraph{Reality of $T_{\mu\nu}[\Theta_I]$.}
With $\Theta_I = i\,\widetilde\Theta_I$:
\begin{equation}
  \partial_\mu\Theta_I^\dagger\partial_\nu\Theta_I
  \;=\; (-i)\,\partial_\mu\widetilde\Theta_I^\dagger
        \cdot (i)\,\partial_\nu\widetilde\Theta_I
  \;=\; \partial_\mu\widetilde\Theta_I^\dagger\,\partial_\nu\widetilde\Theta_I,
\end{equation}
which is \emph{real} (since $\widetilde\Theta_I$ is real-valued).
Similarly $(\partial\Theta_I)^2 = (\partial\widetilde\Theta_I)^2$
and $V(\Theta_I^\dagger\Theta_I) = V(\widetilde\Theta_I^\dagger\widetilde\Theta_I)$
are real.  Therefore:

\begin{theorem}[Reality of imaginary-sector energy-momentum]
\label{thm:tmunu_real}
The energy-momentum tensor of an imaginary-sector object
$\Theta_I = i\,\widetilde\Theta_I$ (with $\widetilde\Theta_I$
a real-valued biquaternion field) is \emph{real}:
\begin{equation}
  T_{\mu\nu}[\Theta_I] \;\in\; \mathbb{R}.
\end{equation}
It acts as a genuine real source in the effective Einstein equations
$G_{\mu\nu} = 8\pi G\,T_{\mu\nu}^{\mathrm{eff}}$ and can, in
principle, curve the real metric $g_{\mu\nu}$.
\end{theorem}

\begin{remark}[Gravitational lensing without direct detection]
Since $T_{\mu\nu}[\Theta_I]$ is real, an imaginary-sector object
\emph{does} produce gravitational effects on real-sector photons.
At the same time, $\Theta_I$ does not couple directly to the
electromagnetic field (which resides in the real sector), so the
imaginary-sector object is electromagnetically invisible.

This is precisely the phenomenological signature of cold dark matter:
gravitationally interacting but electromagnetically silent.
\end{remark}

\paragraph{Cross-term $T_{\mu\nu}^{\mathrm{cross}}$.}
From Propositions~\ref{prop:kin_cross} and \ref{prop:pot_cross},
the first non-vanishing real cross-term in $T_{\mu\nu}$ appears at
order $|\Theta_R||\Theta_I|^2$ (from the second-order potential
correction) and $|\partial\Theta_R||\partial\Theta_I|^2$ (from
fourth-order kinetic corrections).  At leading order in $\Theta_I$,
the sectors are gravitationally \emph{independent}:
\begin{equation}
  T_{\mu\nu}^{\mathrm{eff}}
  \;\approx\; T_{\mu\nu}[\Theta_R] + T_{\mu\nu}[\Theta_I]
  + O(|\Theta_I|^2).
\end{equation}

%----------------------------------------------------------------------
\section{Sector Coupling via the Metric Back-Reaction}
%----------------------------------------------------------------------

Even when $T_{\mu\nu}^{\mathrm{cross}} = 0$ at leading order,
the two sectors are coupled \emph{gravitationally}: $T_{\mu\nu}[\Theta_I]$
sources $g_{\mu\nu}$, and $g_{\mu\nu}$ in turn governs the
dynamics of $\Theta_R$ through the covariant derivative.

This indirect coupling means:
\begin{itemize}
  \item Real-sector objects are deflected by the curvature generated
    by the imaginary sector (gravitational lensing).
  \item There is no direct production or annihilation of imaginary-sector
    quanta in real-sector collisions (to leading order in $\Theta_I$).
  \item The imaginary sector is \emph{dark}: it does not emit, absorb,
    or scatter electromagnetic radiation.
\end{itemize}

%----------------------------------------------------------------------
\section{No-Go Result: Complete Sector Decoupling}
%----------------------------------------------------------------------

\begin{proposition}[No complete decoupling from gravity]
\label{prop:no_decoupling}
An imaginary-sector object $\Theta_I$ with non-zero $\widetilde\Theta_I$
always generates a real, non-zero $T_{\mu\nu}[\Theta_I]$ (by
Theorem~\ref{thm:tmunu_real}) and therefore always sources the
real gravitational field.  Complete decoupling from the real sector
would require $T_{\mu\nu}[\Theta_I] = 0$ everywhere, which (for a
non-trivial $\widetilde\Theta_I$) would demand $V'' = 0$ and all
kinetic terms vanish — a fine-tuned and physically trivial configuration.
\end{proposition}

This is the key result of ST-2: imaginary-sector objects cannot be
completely dark.  They inevitably interact gravitationally with the
real sector.  The two scenarios are:
\begin{enumerate}
  \item \textbf{Partial invisibility (generic case)}: $\Theta_I$
    is electromagnetically invisible but gravitationally detectable
    via lensing, frame dragging, and large-scale structure formation.
  \item \textbf{Complete invisibility (fine-tuned)}: $T_{\mu\nu}[\Theta_I] = 0$
    requires a vanishing or specially tuned potential and kinetic term,
    leaving $\Theta_I$ with no observable effect at all —
    a physically trivial and undetectable configuration.
\end{enumerate}

%----------------------------------------------------------------------
\section{Dark Matter Analogy}
%----------------------------------------------------------------------

The partial-invisibility scenario (case 1 above) reproduces the
key phenomenological properties of cold dark matter (CDM):
\begin{itemize}
  \item Non-luminous: no electromagnetic coupling at leading order.
  \item Gravitationally active: sources Einstein equations through
    $T_{\mu\nu}[\Theta_I]$.
  \item Non-relativistic (in principle): the imaginary-sector potential
    $V(\widetilde\Theta_I^\dagger\widetilde\Theta_I)$ can support
    stable bound configurations analogous to CDM halos.
\end{itemize}

This suggests that UBT may provide a \emph{geometric origin} for
dark matter: the imaginary sector of the biquaternionic field is
naturally dark (electromagnetically) yet gravitationally active,
without requiring the introduction of a new particle species.

\begin{conjecture}[Imaginary sector as geometric dark matter]
\label{conj:dark_matter}
The imaginary-sector component of the UBT $\Theta$-field provides a
natural dark-matter candidate with the following properties:
(i) no electromagnetic coupling at leading order,
(ii) real energy-momentum tensor sourcing GR curvature,
(iii) stability governed by the imaginary-sector potential
     (see \texttt{st3\_stability.tex}).
The relic abundance and mass spectrum would be determined by the
UBT potential $V$ and initial conditions in the biquaternionic field.
\end{conjecture}

Conjecture~\ref{conj:dark_matter} is highly speculative and would
require: (a) identification of the imaginary-sector mass scale,
(b) derivation of the equation of state for $\Theta_I$,
(c) consistency with Big Bang nucleosynthesis and CMB constraints.

%----------------------------------------------------------------------
\section{Summary}
%----------------------------------------------------------------------
\begin{enumerate}
  \item Cross-terms between $\Theta_R$ and $\Theta_I$ in the kinetic
    action are purely imaginary at first order and do not contribute
    to $\operatorname{Re}(S)$.  They appear at second order in $\Theta_I$
    through the potential.
  \item The energy-momentum tensor $T_{\mu\nu}[\Theta_I]$ is real
    (Theorem~\ref{thm:tmunu_real}), so imaginary-sector objects
    gravitate and can produce lensing signatures.
  \item Complete sector decoupling requires fine-tuning that removes
    all dynamical content from $\Theta_I$ — physically trivial.
  \item The generic case is \emph{partial invisibility}:
    $\Theta_I$ is electromagnetically invisible but gravitationally
    detectable, analogous to CDM.  This is the key result.
  \item UBT may provide a geometric origin for dark matter via the
    imaginary sector.
\end{enumerate}

\medskip
\noindent\textbf{All statements in this document are speculative.
No claim constitutes an established result of the canonical UBT framework.}

%----------------------------------------------------------------------
\bibliographystyle{plain}
\begin{thebibliography}{9}
\bibitem{UBT_core}
  D.~Jaroš,
  \emph{Unified Biquaternion Theory: core field equations},
  Repository document \texttt{consolidation\_project/ubt\_core\_main.tex} (2025).
\bibitem{ST1}
  D.~Jaroš,
  \emph{ST-1: Imaginary Sector Objects in UBT — Formal Definition},
  Repository document
  \texttt{speculative\_extensions/invisibility/st1\_imaginary\_object\_definition.tex}
  (2026).
\bibitem{Bertone2005}
  G.~Bertone, D.~Hooper and J.~Silk,
  \emph{Particle dark matter: Evidence, candidates and constraints},
  Phys.\ Rep.\ \textbf{405}, 279--390 (2005).
\bibitem{CTC_conditions}
  D.~Jaroš,
  \emph{ST-3: Closed Timelike Curve Conditions in the Biquaternionic Metric},
  Repository document \texttt{speculative\_extensions/causality/ctc\_conditions.tex}
  (2026).
\end{thebibliography}

\end{document}
